\section{Ekonomija pažnje i nadzorni kapitalizam}

Digitalne platforme ne funkcionišu samo kao prostori komunikacije, zabave i informisanja. One su i ekonomski sistemi u kojima pažnja korisnika postaje jedan od osnovnih resursa. U takvom modelu nije dovoljno da korisnik samo pristupi platformi; važno je da ostane, reaguje, vraća se i ostavlja za sobom podatke koji mogu biti dalje analizirani. Pažnja zato nije samo psihološka sposobnost pojedinca, već i vrednost oko koje se organizuje čitav digitalni poslovni model.

Ekonomija pažnje polazi od činjenice da je ljudska pažnja ograničena. Korisnik ne može istovremeno gledati sve sadržaje i reagovati na sve stimuluse. Upravo zato se platforme takmiče za njegove minute, poglede, klikove i povratke. Ono što uspe da privuče pažnju dobija mogućnost da bude viđeno, monetizovano, preporučeno i dalje algoritamski pojačano.

Međutim, pažnja sama po sebi nije dovoljna za razumevanje savremenih platformi. Njena ekonomska vrednost posebno raste onda kada se poveže sa podacima. Korisnik ne proizvodi vrednost samo time što gleda sadržaj, već i time što ostavlja tragove o sopstvenom ponašanju: šta otvara, koliko dugo se zadržava, šta preskače, na šta reaguje, kada se vraća i u kom trenutku odustaje. Ti tragovi se ne koriste samo za opisivanje ponašanja, već i za njegovo predviđanje. Platforma tako ne posmatra korisnika samo kao publiku, već kao izvor podataka iz kojih se mogu izvoditi obrasci budućeg ponašanja.

Šošana Zubof (Shoshana Zuboff) ovaj model opisuje pojmom nadzornog kapitalizma. U njenom tumačenju, ljudsko iskustvo postaje sirovina za prikupljanje podataka, a ti podaci se zatim koriste za pravljenje prediktivnih proizvoda: procena o tome šta će korisnik verovatno uraditi, pogledati, poželeti, kupiti ili kliknuti \cite{zuboff2015bigother}. Ključna promena nije samo u tome što platforme znaju mnogo o korisniku, već u tome što to znanje nije prvenstveno namenjeno njemu. Ono se koristi unutar tržišta u kome su predviđanja korisničkog ponašanja ekonomski vredna.

U tom smislu, korisnik se nalazi u neobičnoj poziciji. Formalno, on koristi besplatnu ili povoljnu uslugu: društvenu mrežu, pretraživač, aplikaciju za video-sadržaj, platformu za komunikaciju. Međutim, cena te usluge ne plaća se samo novcem. Plaća se pažnjom, podacima i sve preciznijim uvidom platforme u njegove navike. Ono što se korisniku prikazuje kao personalizovano iskustvo istovremeno je i mehanizam prikupljanja novih informacija o njemu. Svaka reakcija popravlja sliku koju sistem ima o korisniku, a ta slika zatim služi za još preciznije oblikovanje narednog iskustva.

Zato personalizacija nije neutralna. Ona može biti korisna kada korisniku olakšava snalaženje u velikoj količini sadržaja, ali u poslovnom modelu zasnovanom na pažnji ona ima i drugu funkciju: povećava verovatnoću da korisnik ostane duže i reaguje češće. Platforma ne mora otvoreno da prisiljava korisnika ni na šta. Dovoljno je da na osnovu podataka o njegovim prethodnim reakcijama, zadržavanju, preskakanju sadržaja i vraćanju uči šta ga zadržava, šta ga uznemirava, šta ga vraća i šta ga čini predvidljivim. U tom trenutku nadzor više ne izgleda kao klasično posmatranje, već kao svakodnevna personalizacija.

Ovakav model menja i način na koji treba posmatrati tehničke metode platformi. Algoritmi preporuke, A/B testiranje, metrike korisničkog angažovanja i psihološki mehanizmi zadržavanja nisu izolovani tehnički trikovi. Oni su delovi istog ekonomskog sistema. Njihova funkcija nije samo da korisniku ponude relevantan sadržaj, već da pažnju učine merljivom, podatke upotrebljivim, ponašanje predvidljivim, a zadržavanje profitabilnim.

Upravo zato se algoritamska manipulacija pažnjom ne može izjednačiti sa „slučajno zlonamernim dizajnom”, niti se njeno postojanje može pobijati argumentom slabosti i nedostatka samokontrole kod pojedinaca. Ono je pitanje odnosa između tehnologije, tržišta i moći. Kada se pažnja pretvori u resurs, ponašanje u podatak, a predviđanje u proizvod, korisničko iskustvo prestaje da bude samo prostor slobodnog izbora. Ono postaje prostor u kome se ekonomski interes platforme neprestano sudara sa autonomijom korisnika.